\section{The exact annihilator of the actual normal
extension}\label{the-exact-annihilator-of-the-actual-normal-extension}

Complete independent derivation. Stable proof locators NEA0--NEA10.

\subsection{NEA0. Construction stage and source
coverage}\label{nea0.-construction-stage-and-source-coverage}

The support remains \(\tau\langle Z_1;\mathrm{no}\ Z_2\rangle\). The
present operations act on the complete arithmetic coefficient receiver
already constructed after whole-history reconstruction. The two branches
retain separate counters. No addition, subtraction, midpoint, metric, or
numerical coordinate is assigned to \(\tau\).

The current \nolinkurl{READ_FIRST_USER_CONSTRUCTION.md,}
correction-precedence table and latest stages of
\nolinkurl{USER_ARGUMENT_RECONSTRUCTION.md,} and PC10--PC11 prerequisite
records were recalled for this continuation. The complete correction
corpus remains the governing record. WU061 concerns the actual lifting
map; WU062 requires checking every operation against the construction;
WU065 requires a concrete attempted next step using the accumulated
findings. Here the input is the actual nonzero normal extension retained
by GMS, and the step is to identify exactly its multiplier annihilator
and then the module generated by the actual connecting morphism.

The direct mathematical prerequisites are
\nolinkurl{CC_GLOBAL_MULTIPLIER_SEPARATION.md,} GMS0--GMS11, completely
derived and independently reviewed at SHA256
\nolinkurl{3164C3AE00FBFBAF17BCE809D48B7D0E56FF8A4C0991FDB374A4B11DD18BEC59};
\nolinkurl{CC_LOCALIZATION_POLYNOMIAL_WEIGHTS.md,} CLP12--CLP13; and
\nolinkurl{CC_ACTUAL_WEIGHT_LIFT_COMPARISON.md,} CW1--CW8. The full
current GMS proof and its actual normal row are the source for the
operations below. The exact source theorem and continuous Schwartz
inverse are the SSI/ESI results received through CW1 and GMS2/GMS8. The
human-source relation is the previously read Connes--Consani §5
construction and the Deligne §3.6 comparison recorded in those notes. No
new human-source reading or new geometric specialization map is claimed
here.

\subsection{NEA1. The actual ring, full ideal, and original
source}\label{nea1.-the-actual-ring-full-ideal-and-original-source}

Use the unital commutative ring \[
M=\{h\in\mathcal O(\mathbb C):\ 
\forall A\ge0\ \exists C_A>0,\ N_A\in\mathbb N,
\ |h(x+it)|\le C_A(1+|t|)^{N_A}\text{ for }|x|\le A\},
\tag{NEA1.1}
\] and the actual Fréchet module \[
\mathcal B=\{F\in\mathcal O(\mathbb C):
b_{A,N}(F)=\sup_{|\Re\lambda|\le A}(1+|\Im\lambda|)^N|F(\lambda)|<\infty
\text{ for all }A,N\ge0\}.
\tag{NEA1.2}
\] Let \(\mathscr Z\) be the original nontrivial-zero divisor of
\(\zeta\), with multiplicities \(m_\rho\). Its translated divisor is \[
\mathscr D_+=\sum_{\rho\in\mathscr Z}m_\rho[\rho+1],\qquad
1<\Re(\rho+1)<2.
\tag{NEA1.3}
\] The support notation counts each distinct zero once, with its
separately retained multiplicity. Define \[
\begin{aligned}
\mathcal I_+&=\{F\in\mathcal B:F^{(j)}(\rho+1)=0
\text{ for }0\le j<m_\rho,\ \rho\in\mathscr Z\},\\
\mathcal Q_+&=\mathcal B/\mathcal I_+,\\
K_+&=\{h\in M:h^{(j)}(\rho+1)=0
\text{ for }0\le j<m_\rho,\ \rho\in\mathscr Z\}.
\end{aligned}
\tag{NEA1.4}
\] Thus \(\mathcal I_+=\mathcal B\cap K_+\); \(K_+\) is an ideal of
\(M\), whereas \(\mathcal I_+\) is the full closed ideal within the
actual source module. They are not identified with each other. Every
multiplier acts continuously on \(\mathcal B\) and its closed invariant
subspaces and quotients, since \[
b_{A,N}(hF)\le C_A b_{A,N+N_A}(F).
\tag{NEA1.5}
\] The derivative product formula preserves all the vanishing orders.

The complete original transform is \[
F_0(s)=\frac{s(s-1)}8\pi^{-s/2}\Gamma(s/2)\zeta(s)
=C_0(s)\zeta(s),
\] \[
F_+(\lambda)=F_0(\lambda-1)
=\frac{(\lambda-1)(\lambda-2)}8
\pi^{-(\lambda-1)/2}\Gamma((\lambda-1)/2)\zeta(\lambda-1).
\tag{NEA1.6}
\] Both are the actual entire source functions from GMS2. In particular
\(F_+\in\mathcal I_+\subset\mathcal B\subset M\), it is nonzero, and its
zero divisor is exactly \(\mathscr D_+\). The original arithmetic is
still \[
\zeta(s)=1+\sum_{n\ge2}n^{-s}=\prod_p(1-p^{-s})^{-1},\qquad
-\frac{\zeta'}{\zeta}(s)=\sum_p\sum_{k\ge1}(\log p)p^{-ks}
\quad(\Re s>1).
\tag{NEA1.7}
\] Thus the unit and every prime-power repetition remain. The
exceptional entire values of the full source factor are \[
F_+(1)=F_+(2)=\frac18,\qquad F_+(0)=\frac\pi{24},
\] \[
\begin{aligned}
F_+(1-2r)
&=\frac{r(2r+1)(-1)^r\pi^r}{2\,r!}\zeta'(-2r)\\
&=F_+(2+2r)
=\frac{(1+2r)(2r)}8\pi^{-(1+2r)/2}
\Gamma((1+2r)/2)\zeta(1+2r)\ne0\quad(r\ge1).
\end{aligned}
\tag{NEA1.8}
\] They retain the pole and Gamma cancellations and the trivial-zero
derivatives; those contributions are not replaced by extra zeros of
\(F_+\).

The object under study is exactly \[
e_+:\quad0\longrightarrow\mathcal I_+\xrightarrow{\iota}\mathcal B
\xrightarrow{q_+}\mathcal Q_+\longrightarrow0,
\qquad
[e_+]\in\operatorname{Ext}^1_M(\mathcal Q_+,\mathcal I_+).
\tag{NEA1.9}
\] We write \(e_+\) for its class when no ambiguity occurs. It is the
translated normal row of CW/CLP, not a separately chosen extension. The
\(M\)-action on the unshifted normal source is \(h(s+1)\); translation
\(VF(\lambda)=F(\lambda-1)\) gives the displayed \(h(\lambda)\)-action
and retains the dilation factor \(a^{s+1}\).

\subsection{NEA2. Every multiplier pushout, with its actual
topology}\label{nea2.-every-multiplier-pushout-with-its-actual-topology}

For \(h\in M\), pushing (NEA1.9) out along multiplication by \(h\) on
\(\mathcal I_+\) gives \[
\mathcal P_h=(\mathcal I_+\oplus\mathcal B)/
\{(hj,-j):j\in\mathcal I_+\},
\quad
i_h(i)=[(i,0)],\quad p_h([(i,b)])=q_+(b).
\tag{NEA2.1}
\] The relation subspace is closed: convergence of \((hj_\nu,-j_\nu)\)
in the product implies convergence of \(j_\nu\) in the closed subspace
\(\mathcal I_+\subset\mathcal B\), and continuity of multiplication
identifies the first limit as multiplication by \(h\) of that limit.
Thus \(\mathcal P_h\) is Fréchet. Each fixed multiplier acts
continuously on it, because it preserves the relation subspace.

The maps form the exact row \[
0\longrightarrow\mathcal I_+\xrightarrow{i_h}\mathcal P_h
\xrightarrow{p_h}\mathcal Q_+\longrightarrow0.
\tag{NEA2.2}
\] Indeed \(i_h\) is injective since a relation with second coordinate
zero has \(j=0\); \(p_h\) is onto; and when \(b\in\mathcal I_+\), the
exact identity \[
[(i,b)]=[(i+hb,0)]
\tag{NEA2.3}
\] identifies its kernel with the image of \(i_h\). These are continuous
maps. The Fréchet open mapping theorem applied to the continuous
bijection \(\mathcal I_+\to\ker p_h\) makes it a topological
isomorphism; applied to the continuous surjection \(p_h\), it gives the
quotient topology. The hypotheses hold because \(\ker p_h\) is closed
and all the spaces are Fréchet. Thus this is a strict extension, not
merely an algebraic presentation endowed with an unspecified topology.
Its Yoneda class is \(h e_+\).

We now prove the splitting criterion directly. If (NEA2.2) has any
algebraic \(M\)-linear section \(s\), define \[
R=i_h^{-1}(1-sp_h):\mathcal P_h\to\mathcal I_+,
\qquad r(b)=R([(0,b)]).
\tag{NEA2.4}
\] Since \((1-sp_h)x\) lies in the kernel, the definition is valid, and
\(R i_h=1\). Formula (NEA2.3) implies \[
r|_{\mathcal I_+}=h|_{\mathcal I_+}.
\tag{NEA2.5}
\] Conversely, if an \(M\)-linear map \(r:\mathcal B\to\mathcal I_+\)
satisfies (NEA2.5), the map \([(i,b)]\mapsto i+r(b)\) is well-defined,
and a section is \[
q\longmapsto[(-r(b),b)]\qquad(q_+(b)=q).
\tag{NEA2.6}
\] Changing \(b\) by \(j\in\mathcal I_+\) changes this pair by
\((-hj,j)\), a relation. Thus \[
h e_+=0\quad\Longleftrightarrow\quad
\text{there is an \(M\)-map }r:\mathcal B\to\mathcal I_+
\text{ extending }h|_{\mathcal I_+}.
\tag{NEA2.7}
\] This proves the specific endomorphism-lifting assertion usually
expressed by the connecting homomorphism in Hom--Ext, without assuming
it as an uncomputed criterion.

\subsection{NEA3. Necessity, sufficiency, and
uniqueness}\label{nea3.-necessity-sufficiency-and-uniqueness}

Suppose \(r\) exists in (NEA2.7). For every \(b\in\mathcal B\),
multiplication by the actual original source \(F_+\in M\) gives \[
F_+r(b)=r(F_+b)=hF_+b.
\tag{NEA3.1}
\] Here \(F_+b\in\mathcal I_+\): \(F_+\) has all the required zeros, and
it is an \(M\)-multiplier of \(\mathcal B\). The middle equality
consequently uses the prescribed restriction (NEA2.5). The nonzero
entire function \(F_+\) can be cancelled from this identity: outside its
discrete zero set the other two factors agree, hence their entire
continuations agree. Thus necessarily \[
r(b)=hb\quad\text{for every }b\in\mathcal B,
\qquad h\mathcal B\subset\mathcal I_+.
\tag{NEA3.2}
\] To read the precise multiplicities from this inclusion, take the
single zero-free source-space test function \[
g(\lambda)=e^{\lambda^2}\in\mathcal B.
\tag{NEA3.3}
\] It belongs to the existing coefficient space because
\(b_{A,N}(g)\le e^{A^2}\sup_t(1+|t|)^Ne^{-t^2}<\infty\). It is a unit in
each local holomorphic ring. Thus \(hg\in\mathcal I_+\) forces
\(\operatorname{ord}_{\rho+1}h\ge m_\rho\) for every actual zero. This
test is a function in the already constructed source space, not an
alternative arithmetic or an alternative zeta function.

Conversely, these full vanishing orders make \(hb\in\mathcal I_+\) for
every \(b\in\mathcal B\) by the derivative product formula. The map
\(r_h(b)=hb\) is continuous by (NEA1.5), so (NEA2.6) gives a continuous
splitting. We have proved \[
\boxed{\operatorname{Ann}_M(e_+)=K_+
=\{h\in M:\operatorname{ord}_{\rho+1}h\ge m_\rho\text{ for every }\rho\in\mathscr Z\}.}
\tag{NEA3.4}
\] The same source-space test shows \[
K_+=\operatorname{Ann}_M(\mathcal Q_+),
\tag{NEA3.5}
\] since annihilation of the quotient is exactly the inclusion
\(h\mathcal B\subset\mathcal I_+\).

For every \(h\in K_+\), the full continuous isomorphism of its pushout
with the split row is \[
\Phi_h([(i,b)])=(i+hb,q_+(b)),\qquad
\Phi_h^{-1}(i_0,q)=[(i_0-hb,b)]\quad(q_+(b)=q).
\tag{NEA3.6}
\] The first coordinate is in the full \(\mathcal I_+\). Both formulas
are independent of representatives by (NEA2.1). Their composites are
identities by substitution. For continuity of the inverse, the map
\((i_0,b)\mapsto[(i_0-hb,b)]\) is continuous and kills
\(\{0\}\oplus\mathcal I_+\); it therefore descends through the product
quotient. The resulting section is \(q\mapsto[(-hb,b)]\).

Every such section is unique, even among algebraic \(M\)-linear
sections. Indeed \[
\operatorname{Hom}_M(\mathcal Q_+,\mathcal I_+)=0,
\qquad \operatorname{Hom}_M(\mathcal Q_+,\mathcal B)=0.
\tag{NEA3.7}
\] To prove both statements, observe that \(F_+\) annihilates
\(\mathcal Q_+\) and acts injectively by multiplication on either
entire-function target. For an \(M\)-map \(v\), the identity
\(F_+v(q)=v(F_+q)=0\) therefore forces \(v=0\). The difference of two
sections of (NEA2.2) is a map into its kernel, so it is zero.
Consequently any algebraic splitting that exists is automatically the
explicit continuous one in (NEA3.6); no continuity of a hypothetical
algebraic section was assumed in the necessity proof.

The actual extension is nonzero: \(1\notin K_+\), since the actual
divisor is nonempty, as established in CW8. Thus (NEA3.4) includes
\(e_+\ne0\). It also proves that no classification of arbitrary
endomorphisms of \(\mathcal B\), no principal-ideal assertion, and no
unproved growth bound for \(h/F_+\) is needed.

\subsection{NEA4. Pullback of every multiplier on the
quotient}\label{nea4.-pullback-of-every-multiplier-on-the-quotient}

For clarity, lifting the quotient endomorphism \(h_{\mathcal Q_+}\) is a
different presentation of the same killed class. Its actual pullback is
\[
\mathcal E_h=\{(b,q)\in\mathcal B\oplus\mathcal Q_+:
q_+(b)=h q\},
\quad i'_h(j)=(j,0),\quad p'_h(b,q)=q.
\tag{NEA4.1}
\] It is a closed subspace of the product, with its Fréchet topology,
and is a strict exact extension by the same open-mapping argument as in
NEA2. There is an explicit continuous \(M\)-linear isomorphism from the
pushout to this pullback: \[
\mathcal P_h\longrightarrow\mathcal E_h,
\qquad [(i,b)]\longmapsto(i+hb,q_+(b)).
\tag{NEA4.2}
\] For arbitrary \(h\), the first coordinate here is in \(\mathcal B\),
not asserted to be in \(\mathcal I_+\). It has the required quotient
\(h q_+(b)\). The relations are killed. The map is onto: if
\((b_0,q)\in\mathcal E_h\), choose a lift \(b\) of \(q\); then
\(i=b_0-hb\in\mathcal I_+\), and \([(i,b)]\) maps to \((b_0,q)\). Its
kernel is exactly the relation subspace, by (NEA2.3). Thus it is a
continuous bijection of Fréchet spaces and its inverse is continuous by
the open mapping theorem. It is the identity on the displayed kernel and
quotient. This proves concretely that pushout and pullback define the
same scalar action \(h e_+\).

A section of (NEA4.1) is equivalent to an \(M\)-map
\(v:\mathcal Q_+\to\mathcal B\) with \(q_+v=h_{\mathcal Q_+}\). By
(NEA3.7) the only possible \(v\) is zero. Hence this lift exists
precisely when \(h\in K_+\), and then the section is \(q\mapsto(0,q)\).
Under (NEA4.2) it is exactly \(q\mapsto[(-hb,b)]\). This does not
identify the nonzero source-return map
\(b\mapsto hb:\mathcal B\to\mathcal I_+\) with the zero
quotient-to-source lift; their domains and purposes are different and
their exact relation is (NEA4.2).

\subsection{NEA5. Full multiplicity probes and exact residual
divisors}\label{nea5.-full-multiplicity-probes-and-exact-residual-divisors}

Fix an actual \(\lambda_0=\rho+1\), put \(z=\lambda-\lambda_0\), and
retain \(m=m_\rho\). The full local source factor is \[
F_+(\lambda_0+z)=z^m b_\rho(z),\qquad
b_\rho(z)=C_0(\rho+z)\frac{\zeta(\rho+z)}{z^m},\qquad
b_\rho(0)=C_0(\rho)\frac{\zeta^{(m)}(\rho)}{m!}\ne0.
\tag{NEA5.1}
\] For example all its Taylor coefficients, without dropping the
completion derivatives, are \[
[z^j]b_\rho(z)
=\sum_{\ell=0}^j
\frac{C_0^{(\ell)}(\rho)}{\ell!}
\frac{\zeta^{(m+j-\ell)}(\rho)}{(m+j-\ell)!}.
\tag{NEA5.2}
\] Here \(C_0(s)=s(s-1)\pi^{-s/2}\Gamma(s/2)/8\) is the full function in
(NEA1.6).

Division by a fixed polynomial preserves \(M\) whenever the numerator
vanishes at every root to the needed order. To verify this, remove fixed
small disks around the finitely many roots; on the complement the
reciprocal polynomial has at most a fixed polynomial bound, and in each
disk the entire quotient is bounded on a compact set. The same assertion
holds for \(\mathcal B\), retaining all its rapid strip seminorms. In
particular \(M/(z^k)\) is exactly \(\mathbb C[z]/z^k\): the Taylor
polynomial supplies every jet, and a zero jet is divisible by \(z^k\)
with quotient in \(M\).

For \(1\le k\le m\), the actual full source gives the \(M\)-linear
injection \[
\iota_{\rho,k}:M/(z^k)\longrightarrow\mathcal Q_+,
\qquad [a]\longmapsto\left[a\frac{F_+}{z^k}\right].
\tag{NEA5.3}
\] The source function \(F_+/z^k\) is entire and in \(\mathcal B\); at
the chosen zero it has exact order \(m-k\), and at every other zero it
retains the required order. Its product with \(a\) lies in
\(\mathcal I_+\) exactly when \(a\) has a zero of order at least \(k\)
at \(\lambda_0\). This proves both well-definedness and injectivity,
with every other zero still retained.

The two free \(M\)-modules in
\(0\to M\xrightarrow{z^k}M\to M/(z^k)\to0\) form a resolution, since
\(M\) is a ring of entire functions and multiplication by \(z^k\) is
injective. Pulling the class \(h e_+\) back along (NEA5.3) gives exactly
\[
\iota_{\rho,k}^*(h e_+)=[hF_+]
\in\operatorname{Ext}^1_M(M/(z^k),\mathcal I_+)
=\mathcal I_+/z^k\mathcal I_+.
\tag{NEA5.4}
\] To check the positive sign, use the lift \(F_+/z^k\) of the generator
in the original extension. Its relation is \(z^k(F_+/z^k)=F_+\); pushing
its kernel by \(h\) changes that relation to \(hF_+\). This is the
positive boundary convention used in CLP12.

There is an exact full-factor coordinate and inverse \[
\mathcal I_+/z^k\mathcal I_+\xrightarrow{\sim}\mathbb C[z]/z^k,
\quad [j]\longmapsto\operatorname{Tay}_{\lambda_0,<k}(j/F_+),
\qquad [p]\longmapsto[F_+p].
\tag{NEA5.5}
\] The quotient \(j/F_+\) in this formula is taken locally at
\(\lambda_0\), where (NEA5.1) is the retained unit comparison. If its
first \(k\) coefficients vanish, then \(j/z^k\) is entire, lies in
\(\mathcal B\), and lies in \(\mathcal I_+\): at the chosen point it
still vanishes to order at least \(m\), and at the other points division
is by a unit. This proves the kernel claim. Multiplication by \(F_+\)
times a polynomial proves surjectivity and the two inverse identities.
Derivative evaluations are continuous on the actual Fréchet ideal. Thus
this finite-dimensional coordinate is continuous as well; no global
claim that \(j/F_+\in M\) is made.

Under (NEA5.5), the boundary (NEA5.4) is the complete \(k\)-jet of
\(h\). In particular the \(k=m\) probe vanishes precisely when
\(\operatorname{ord}_{\lambda_0}h\ge m\). The eigenvector probe \(k=1\)
detects only \(h(\lambda_0)\) and would miss higher multiplicities. If
\(v=\operatorname{ord}_{\lambda_0}h\), the local class has exact
\(z\)-nilpotence exponent \(\max(k-v,0)\), with exponent zero meaning
the class is zero. Indeed its first nonzero Taylor term is a nonzero
coefficient times \(z^v\), and multiplication by \(z\) raises this order
by exactly one until it reaches \(k\).

For the complete global class the exact residual annihilator is
consequently \[
\operatorname{Ann}_M(h e_+)=(K_+:h)
=\{a\in M:\operatorname{ord}_{\rho+1}a\ge
\max(m_\rho-\operatorname{ord}_{\rho+1}h,0)\text{ for every }\rho\}.
\tag{NEA5.6}
\] This follows also directly by applying (NEA3.4) to \(ah\). When
\(h=0\), all residual multiplicities are zero and the annihilator is
\(M\). This classifies the effect of every multiplier on the actual
extension, retaining its entire divisor rather than only finitely many
tests.

\subsection{NEA6. The actual cyclic module of extension
classes}\label{nea6.-the-actual-cyclic-module-of-extension-classes}

Define the actual image \[
\mathscr N_+=M e_+\subset\operatorname{Ext}^1_M(\mathcal Q_+,\mathcal I_+).
\tag{NEA6.1}
\] It is the image of the connecting homomorphism on the specified
multiplier endomorphisms of \(\mathcal I_+\). No claim that these are
all endomorphisms is needed. The exact map and its kernel give \[
M/K_+\xrightarrow{\sim}\mathscr N_+,
\qquad [h]\longmapsto h e_+.
\tag{NEA6.2}
\] The isomorphism is canonical for the actual displayed extension, with
\([1]\) mapped to \(e_+\). It is an algebraic \(M\)-module statement. We
assign no unproved locally convex quotient topology to \(M/K_+\) or to
the ambient Ext group.

The full jet map gives a faithful concrete description of this module:
\[
M/K_+\longrightarrow\prod_{\rho\in\mathscr Z}\mathbb C[z]/z^{m_\rho},
\quad [h]\longmapsto
\left(\sum_{j=0}^{m_\rho-1}\frac{h^{(j)}(\rho+1)}{j!}z^j\right)_\rho.
\tag{NEA6.3}
\] It is injective by definition of \(K_+\). Its image is the image of
the actual strip-polynomial multiplier functions. It is not declared to
be the unrestricted product. The projection onto any specified finite
set of these jet factors is surjective: the usual finite polynomial
remainder construction follows from Bézout for the pairwise coprime
polynomials \((\lambda-\rho-1)^{m_\rho}\), and the resulting polynomial
belongs to \(M\). This finite assertion neither supplies arbitrary
infinite interpolation nor assumes convergence of finite spectral
expansions.

There is also the exact injective \(M\)-map from the original normal
quotient \[
\mathcal Q_+=\mathcal B/\mathcal I_+\longrightarrow M/K_+,
\qquad [F]_{\mathcal I_+}\longmapsto[F]_{K_+}.
\tag{NEA6.4}
\] It is induced by \(\mathcal B\subset M\), and injectivity is
precisely \(\mathcal I_+=\mathcal B\cap K_+\). In terms of (NEA6.2), it
sends \([F]\) to \(F e_+\). We do not identify the two full quotient
modules or assert surjectivity of (NEA6.4).

The ideal \(K_+\), hence the annihilator of the single actual extension
class, recovers the complete translated zero divisor: its common zero
set is exactly the zero set of \(F_+\), since \(F_+\in K_+\), and at
\(\lambda_0=\rho+1\), \[
\min_{0\ne h\in K_+}\operatorname{ord}_{\lambda_0}h=m_\rho.
\tag{NEA6.5}
\] Every member has at least that order, and the original full factor
\(F_+\) attains it. This is a precise recovery of the divisor and its
multiplicities from the annihilator, not a claim about where the
original zeros lie within their strip.

\subsection{NEA7. The exact full-module generator
resolvent}\label{nea7.-the-exact-full-module-generator-resolvent}

Write \(T[h]=[\lambda h]\) on \(M/K_+\), equivalently postcomposition by
the normal generator on \(\mathscr N_+\). For a scalar
\(\mu\notin\operatorname{Supp}\mathscr D_+\), define \[
\mathcal R_\mu h(\lambda)=
\frac{h(\lambda)-F_+(\lambda)h(\mu)/F_+(\mu)}{\lambda-\mu}.
\tag{NEA7.1}
\] The denominator of the retained full scalar is nonzero. The numerator
has a zero at \(\mu\), so the quotient is entire. It belongs to \(M\):
outside a fixed disk around \(\mu\) this follows from the strip bounds
for \(h\) and \(F_+\), while on the disk the entire quotient is bounded
on a compact set. If \(h\in K_+\), both numerator terms vanish to every
required translated order, and the denominator is a unit at each
translated zero. Thus \(\mathcal R_\mu\) preserves \(K_+\) and induces
an operator on the whole quotient.

Multiplication gives the two exact identities \[
(T-\mu)\mathcal R_\mu[h]=[h],\qquad
\mathcal R_\mu(T-\mu)[h]=[h].
\tag{NEA7.2}
\] For the first, the retained difference from \(h\) is
\(F_+h(\mu)/F_+(\mu)\in K_+\); for the second, the evaluated numerator
\(((\lambda-\mu)h)(\mu)\) is zero. The inverse is \(M\)-linear on the
quotient: \(T-\mu\) is \(M\)-linear and a bijection, so its inverse is
\(M\)-linear. Formula (NEA7.1) already specifies that inverse on each
representative, without a finite-block-density argument.

At an actual \(\lambda_0=\rho+1\), the nonzero class \[
\left[\frac{F_+}{\lambda-\lambda_0}\right]_{K_+}
\tag{NEA7.3}
\] is an eigenvector of \(T\) with eigenvalue \(\lambda_0\): it has
order \(m_\rho-1\) there and is killed by \(T-\lambda_0\). Hence
\(T-\lambda_0\) is not injective. We have therefore computed the entire
algebraic resolvent set and its complement: \[
\{\mu\in\mathbb C:T-\mu\text{ is not bijective on }\mathscr N_+\}
=\operatorname{Supp}\mathscr D_+=\{\rho+1:\rho\in\mathscr Z\}.
\tag{NEA7.4}
\] This statement specifies algebraic bijectivity; it does not invoke
the spectrum of an operator on an unassigned Banach or Fréchet topology.

The corresponding complete local chains are retained as actual
submodules. At \(\lambda_0\), the map \[
M/((\lambda-\lambda_0)^m)\longrightarrow M/K_+,
\quad[a]\longmapsto\left[a\frac{F_+}{(\lambda-\lambda_0)^m}\right]
\tag{NEA7.5}
\] is injective by the same full-order argument as (NEA5.3). On this
block, multiplication by \(a_0^\lambda\), for recovered \(a_0>0\), is
exactly \[
a_0^{\rho+1}\sum_{j=0}^{m-1}\frac{(\log a_0)^j}{j!}M_z^j.
\tag{NEA7.6}
\] The map includes the entire factor in (NEA5.1); no derivative
coefficients are removed by this expression for the action.

Every unital complex-algebra character of \(M/K_+\) is evaluation at
exactly one translated zero. To prove this, let \(\chi\) be such a
character and put \(\mu=\chi([\lambda])\). For every \(h\in M\), the
function \((h(\lambda)-h(\mu))/(\lambda-\mu)\) is entire and belongs to
\(M\), by the division estimate used above. Applying \(\chi\) gives
\(\chi([h])=h(\mu)\). Since \(F_+\in K_+\), it follows that
\(F_+(\mu)=0\), so \(\mu\in\operatorname{Supp}\mathscr D_+\). Conversely
evaluation at any such point annihilates \(K_+\) and defines a
character. Thus their actual dilation characters are \(a_0^{\rho+1}\).
For \(a_0>1\), their logarithmic-modulus values are \[
\frac{2\log|a_0^{\rho+1}|}{\log a_0}
=2\Re\rho+2\in(2,4).
\tag{NEA7.7}
\] These are the scalar values for the postcomposition module action
just defined. They are not weights of a conjugation action on
equivariant Hom; NEA10 retains that distinction.

\subsection{NEA8. Separation of the actual class module from the
original
quotient}\label{nea8.-separation-of-the-actual-class-module-from-the-original-quotient}

The existing global separator is
\(c(\lambda)=E_+(\lambda)=1-E(\lambda-1)\). GMS3 proves that
\(c\in K_+\), with full orders, and that it acts as the identity on the
complete original quotient \(\mathcal Q=\mathcal B/\mathcal I\).
Consequently \[
c\mathscr N_+=0,\qquad c\mathcal Q=\mathcal Q
\text{ with }c_{\mathcal Q}=1.
\tag{NEA8.1}
\] The second equality alone is not the reason; its stated identity
action is the property used.

Apply the explicit bar-resolution homotopy of GMS4 to these actual
modules. If \(P\to\mathcal Q\) is a projective resolution,
multiplication by \(1-c\) lifts zero on its augmentation, and there is
\(h:P\to P[-1]\) with \(dh+hd=1-c\). The Hom homotopy
\(H^n(f)=(-1)^nfh\) satisfies \[
DH+HD=1-c_{\mathscr N_+}=1
\quad\text{on }\operatorname{Hom}_M(P,\mathscr N_+).
\tag{NEA8.2}
\] For the reverse direction, resolve \(\mathscr N_+\), lift its zero
action of \(c\), and use the same Hom formula; the target action of
\(c\) on \(\mathcal Q\) is one, giving the reverse contraction. Hence \[
\boxed{\operatorname{RHom}_M(\mathcal Q,\mathscr N_+)\simeq0,
\qquad\operatorname{RHom}_M(\mathscr N_+,\mathcal Q)\simeq0.}
\tag{NEA8.3}
\] The complete bar insertion formulas in GMS4.5--GMS4.9 apply without
any finiteness condition on \(\mathscr N_+\). This is separation of a
module generated by the actual extension class. It does not assign the
translated-zero spectrum to the entire kernel \(\mathcal I_+\), whose
other quotient characters remain as calculated in CW6.

\subsection{NEA9. Original-source maps for every annihilating
multiplier}\label{nea9.-original-source-maps-for-every-annihilating-multiplier}

For every \(h\in K_+\), the unique source-return map in (NEA3.2)
becomes, in the original normal variable, \[
V^{-1}r_hV(F)(s)=h(s+1)F(s).
\tag{NEA9.1}
\] Its exact source representative is \[
\mathfrak r_h(a)(u)=\Theta^{-1}\bigl[h(s+1)\Theta a(s)\bigr](u)
=\frac{u^{-1/2}}\pi\int_{\mathbb R}
h(3/2+it)\Theta a(1/2+it)u^{-it}\,dt.
\tag{NEA9.2}
\] The shift \(3/2+it\) is required by the actual normal action. The
multiplier bound, rapid strip decay, and the full inverse Mellin
estimates from GMS2 prove absolute differentiated convergence and
continuity \(A(-1)\to J(-1)\). It lands in \(J\) because \(h(s+1)\) has
the required original-zero orders and \(\Theta J=\mathcal I\). Composing
with the exact continuous inverse \(H=\Sigma^{-1}:J\to S\) gives \[
f_{h,a}=H\mathfrak r_h(a),\qquad
\Sigma f_{h,a}=\mathfrak r_h(a),\qquad
f_{h,a}(0)=0,\quad\int_{\mathbb R}f_{h,a}=0.
\tag{NEA9.3}
\] All sums, factors \(2\), \(1/2\), \(u^{-1/2}/\pi\), and endpoint
conditions are the original ones from GMS2/GMS8. For \(h=c\), formula
(NEA9.2) reduces exactly to multiplication by \(1-E(s)\), because
\(c(s+1)=1-E(s)\). The classification (NEA3.4) says precisely that
these, and only these multipliers, admit the displayed lift into the
full source kernel.

\subsection{NEA10. The actual geometric connecting morphism and its
class
module}\label{nea10.-the-actual-geometric-connecting-morphism-and-its-class-module}

The simultaneous complexes of the actual sheaf row, with their
constructed \(M\)-actions, are \[
D_J=(W\xrightarrow{r_+-r_-}J\xrightarrow0J(-1)),\quad
D_F=(W\xrightarrow{r_+-r_-}A\xrightarrow0A(-1)),\quad
D_G=Q[-1]\oplus Q(-1)[-2].
\tag{NEA10.1}
\] The low generator is \(L\), the normal generator is \(L+1\), and GMS9
retains the endpoint characters \((h(0),h(1),h(1),h(0))\). All
identifications below use \(\Theta\) on low coefficients and \(V\Theta\)
on normal coefficients. The explicit symmetric source contraction in
CLP12/GMS9 gives \[
D_J\simeq Z[0]\oplus\mathcal I_+[-2],\quad
D_F\simeq Z[0]\oplus\mathcal Q[-1]\oplus\mathcal B[-2],\quad
D_G\simeq\mathcal Q[-1]\oplus\mathcal Q_+[-2].
\tag{NEA10.2}
\] These are the compatible \(M\)-linear quasi-isomorphisms of the
actual row, and retain the entire ideal. Its connecting morphism is \[
\delta:D_G\longrightarrow D_J[1],\qquad
\delta|_{\mathcal Q[-1]}=0,\qquad
\delta|_{\mathcal Q_+[-2]}=e_+[-2]
\text{ into }\mathcal I_+[-1].
\tag{NEA10.3}
\] Let \(M\) act on this derived Hom group by postcomposition on
\(D_J[1]\); since \(\delta\) is \(M\)-linear, this equals precomposition
on its source. Inclusion and projection of the normal summands in
(NEA10.2) prove both directions of \[
h\delta=0\quad\Longleftrightarrow\quad h e_+=0.
\tag{NEA10.4}
\] Therefore the actual cyclic boundary-image module is \[
M\delta\subset\operatorname{Hom}_{D(M)}(D_G,D_J[1]),\qquad
M/K_+\xrightarrow{\sim}M\delta,\quad[h]\mapsto h\delta,
\] \[
M\delta\xrightarrow{\sim}\mathscr N_+,
\qquad h\delta\longmapsto h e_+.
\tag{NEA10.5}
\] This identifies the class receiver directly inside the actual
localization triangle; it does not replace its full source or kernel.
Its exact annihilator, resolvent, and full translated divisor are
(NEA3.4), (NEA7.1)--(NEA7.4), and (NEA6.5).

In particular \(c\delta=0\), whereas \(c\) acts as the identity on
\(\mathcal Q[k]\) for every integer \(k\). For every derived
\(M\)-morphism \(f:\mathcal Q[k]\to D_G\), naturality of the central
action gives the actual lifting equation \[
\delta f
=\delta f c_{\mathcal Q[k]}
=\delta c_{D_G}f
=c_{D_J[1]}\delta f
=0.
\tag{NEA10.6}
\] It follows that \(f\) lifts through \(D_F\to D_G\). To display the
lift with the actual decompositions, the normal component of \(f\) is
zero in the derived category by GMS4.4. Write its remaining component as
\(f_1:\mathcal Q[k]\to\mathcal Q[-1]\). Then \((0,f_1,0)\) in the middle
decomposition of (NEA10.2) is a lift, transported back through its
specified quasi-isomorphism. This proves existence; it does not claim
uniqueness of all lifts through \(D_F\), whose differences can come from
\(D_J\).

The action used in (NEA10.5)--(NEA10.6) must not be confused with
conjugation by dilation. Postcomposition by the multiplier \(a^\lambda\)
produces the module action whose scalar values are (NEA7.6)--(NEA7.7).
The conjugation action on a derived morphism is instead \[
\delta\longmapsto T_{a,D_J[1]}\,\delta\,T_{a,D_G}^{-1}.
\tag{NEA10.7}
\] It fixes \(\delta\), because \(\delta\) is equivariant. No
translated-zero weight is assigned to this conjugation action. Likewise
no identification of \(M\delta\) with Deligne's geometric obstruction
quotient is asserted without its comparison map. The proved conclusion
is the explicit annihilator, spectrum and separation of the module
generated by the actual localization class, and the resulting actual
lift (NEA10.6), with all original kernels, factors, multiplicities and
degrees retained.
